\section{Sistemas coordenados}

El estudio del cálculo multivariable requiere entender de forma clara los espacios geométricos donde se desarrollan las funciones de varias variables. A continuación, se desarrollan los conceptos del plano y el espacio euclidiano de forma detallada.

\subsection{El Plano Cartesiano ($\mathbb{R}^2$)}

El sistema de coordenadas rectangulares en el plano se construye mediante dos rectas reales perpendiculares que se intersectan en un punto llamado \textbf{origen}.

\begin{mydefinition}{El Plano $\mathbb{R}^2$}
El conjunto $\mathbb{R}^2$ es el conjunto de todos los pares ordenados de números reales. Matemáticamente se define como el producto cartesiano $\mathbb{R} \times \mathbb{R}$:
$$ \mathbb{R}^2 = \{ (x, y) \mid x \in \mathbb{R}, y \in \mathbb{R} \} $$
A cada par ordenado $(x, y)$ se le denomina \textbf{punto} en el plano.
\end{mydefinition}

Los elementos del par ordenado tienen los siguientes nombres:
\begin{itemize}
    \item $x$: Coordenada $x$ o \textbf{abscisa}. Representa la distancia dirigida desde el eje vertical.
    \item $y$: Coordenada $y$ u \textbf{ordenada}. Representa la distancia dirigida desde el eje horizontal.
\end{itemize}

Los ejes coordenados dividen al plano en cuatro regiones denominadas \textbf{cuadrantes}, numerados en sentido antihorario comenzando por el superior derecho (donde ambas semiejes son positivas).

\begin{center}
\shorthandoff{>}
\begin{tikzpicture}[scale=1.2]
    % Ejes
    \draw[->, thick] (-3,0) -- (3,0) node[right] {$x$ (Eje de abscisas)};
    \draw[->, thick] (0,-3) -- (0,3) node[above] {$y$ (Eje de ordenadas)};
    
    % Punto P
    \coordinate (P) at (2, 1.5);
    \draw[dashed, color=gray] (2,0) node[below, black] {$a$} -- (P);
    \draw[dashed, color=gray] (0,1.5) node[left, black] {$b$} -- (P);
    \filldraw[blue] (P) circle (2pt) node[anchor=south west] {$P(a,b)$};
    
    % Etiquetas de cuadrantes
    \node[gray] at (2,2) {I $(+,+)$};
    \node[gray] at (-2,2) {II $(-,+)$};
    \node[gray] at (-2,-2) {III $(-,-)$};
    \node[gray] at (2,-2) {IV $(+,-)$};
    
    % Origen
    \node[below left] at (0,0) {$O(0,0)$};
\end{tikzpicture}
\end{center}

\subsection{El Espacio Euclidiano ($\mathbb{R}^3$)}

Para representar puntos en el espacio tridimensional, extendemos el concepto anterior añadiendo un tercer eje perpendicular a los otros dos.

\begin{mydefinition}{El Espacio $\mathbb{R}^3$}
El conjunto $\mathbb{R}^3$ es el conjunto de todas las ternas ordenadas de números reales:
$$ \mathbb{R}^3 = \{ (x, y, z) \mid x, y, z \in \mathbb{R} \} $$
\end{mydefinition}

\subsubsection{Orientación y la Regla de la Mano Derecha}
En ingeniería y física, es estándar utilizar un sistema de coordenadas \textit{dextrógiro} (mano derecha). 
\begin{mynote}{Regla de la Mano Derecha}
Si alineas los dedos de tu mano derecha en la dirección del eje $x$ positivo y los cierras hacia el eje $y$ positivo, tu pulgar apuntará en la dirección del eje $z$ positivo.
\end{mynote}

En este sistema:
\begin{itemize}
    \item El eje $z$ suele representar la altura o la cota vertical.
    \item Los tres ejes definen tres \textbf{planos coordenados}:
    \begin{itemize}
        \item \textbf{Plano $xy$}: Formado por los ejes $x$ e $y$ (donde $z=0$).
        \item \textbf{Plano $xz$}: Formado por los ejes $x$ y $z$ (donde $y=0$).
        \item \textbf{Plano $yz$}: Formado por los ejes $y$ y $z$ (donde $x=0$).
    \end{itemize}
\end{itemize}

Estos planos dividen el espacio en ocho regiones llamadas \textbf{octantes}. El primer octante es aquel donde $x, y, z > 0$.

\subsubsection{Localización de un punto en el espacio}

Para localizar un punto $P(a, b, c)$ en el espacio:
1. Avanzamos $a$ unidades sobre el eje $x$.
2. Nos movemos $b$ unidades paralelamente al eje $y$.
3. Subimos (o bajamos) $c$ unidades paralelamente al eje $z$.

\begin{center}
\shorthandoff{>}
\begin{tikzpicture}[x={(-0.5cm,-0.5cm)}, y={(1cm,0cm)}, z={(0cm,1cm)}, scale=1.5]
    % Ejes (parte negativa punteada)
    \draw[dashed, gray] (0,0,0) -- (-3,0,0); % x back
    \draw[dashed, gray] (0,0,0) -- (0,-2,0); % y back
    \draw[dashed, gray] (0,0,0) -- (0,0,-1); % z back
    
    % Ejes positivos
    \draw[->, thick] (0,0,0) -- (3,0,0) node[below left] {$x$};
    \draw[->, thick] (0,0,0) -- (0,4,0) node[right] {$y$};
    \draw[->, thick] (0,0,0) -- (0,0,3) node[above] {$z$};
    
    % Coordenadas del punto
    \def\x{2}
    \def\y{3}
    \def\z{2}
    
    % Construcción de la caja
    \draw[dashed, blue] (0,0,0) -- (\x,0,0) node[left, black] {$a$};
    \draw[dashed, blue] (\x,0,0) -- (\x,\y,0);
    \draw[dashed, blue] (0,\y,0) -- (\x,\y,0);
    \draw[dashed, blue] (0,0,0) -- (0,\y,0) node[below, black] {$b$};
    
    \draw[dashed, blue] (\x,\y,0) -- (\x,\y,\z); % Línea vertical al punto
    
    \draw[dashed, blue] (0,0,\z) -- (\x,0,\z);
    \draw[dashed, blue] (0,0,\z) -- (0,\y,\z);
    \draw[dashed, blue] (\x,0,\z) -- (\x,\y,\z);
    \draw[dashed, blue] (0,\y,\z) -- (\x,\y,\z);
    \draw[dashed, blue] (0,0,0) -- (0,0,\z) node[left, black] {$c$};
    \draw[dashed, blue] (\x,0,0) -- (\x,0,\z);
    \draw[dashed, blue] (0,\y,0) -- (0,\y,\z);

    % Punto y planos etiqueta
    \filldraw[red] (\x,\y,\z) circle (2pt) node[anchor=south west] {$P(a,b,c)$};
\end{tikzpicture}
\end{center}